% ============================================================================
%  MATEMÁTICAS OPERATIVAS — Sesión 5
%  Semana 5: 7 al 12 de septiembre de 2026
%  1.5 Razones y proporciones · Repaso integrador de la Unidad 1
%  Duración prevista: 3 horas
% ============================================================================
\documentclass[aspectratio=169,10pt]{beamer}

\usepackage[utf8]{inputenc}
\IfFileExists{spanish.ldf}{\usepackage[spanish,es-noshorthands]{babel}}{}
\usepackage{amsmath,amssymb}
\usepackage{booktabs}
\usepackage{tikz}

\usetheme{Uniminuto}

\uniasignatura{Matemáticas Operativas}
\uniprofesor{Andrés Felipe Puerta Velez}
\unicorreo{andres.puerta-v@uniminuto.edu.co}
\uninrc{CBASBL011}
\unisemana{Sesión 5 — Semana 5 · 7 al 12 de septiembre de 2026}

\newcommand{\listatiny}{\setbeamerfont{itemize/enumerate body}{size=\tiny}\tiny}

\begin{document}

\portada

\begin{frame}{Orden del día}
  \begin{enumerate}
    \item Revisión de los ejercicios propuestos de la sesión 4
    \item \azul{Tema 1.5: razones, proporciones y regla de tres}
    \item Porcentajes: los cuatro tipos de pregunta
    \item Repartos proporcionales
    \item \azul{Repaso integrador de la Unidad 1}
    \item Simulacro del Examen 1 y recomendaciones
  \end{enumerate}

  \vspace{0.25cm}
  {\footnotesize\color{uniGris} Sesión de 3 horas. Cerramos la unidad: el
  \destacar{Examen 1 (35\%)} es la próxima semana (14--19 de septiembre) en
  Aulas Virtuales, presencial en el aula.}
\end{frame}

% ===========================================================================
\seccion{1.5 Razones y proporciones}
% ===========================================================================

\begin{frame}{Razón y proporción}
  \scriptsize
  \begin{columns}[T,onlytextwidth]
    \begin{column}{0.48\textwidth}
      \begin{block}{Razón}
        \tiny
        Es el \azul{cociente} de dos cantidades: $\dfrac{a}{b}$, también
        escrito $a:b$. Compara \destacar{por división}, no por resta.

        \vspace{0.05cm}
        Si en un grupo hay $18$ mujeres y $12$ hombres, la razón es
        $\tfrac{18}{12} = \tfrac{3}{2}$: \azul{3 mujeres por cada 2 hombres}.
      \end{block}
    \end{column}
    \begin{column}{0.48\textwidth}
      \begin{block}{Proporción}
        \tiny
        Es la \azul{igualdad} de dos razones:
        \[ \frac{a}{b} = \frac{c}{d} \]
        Y la propiedad que lo resuelve todo, el
        \destacar{producto cruzado}:
        \[ a\cdot d = b\cdot c \]
      \end{block}
    \end{column}
  \end{columns}

  \vspace{0.12cm}
  \begin{exampleblock}{Ejemplo}
    \tiny
    Halle $x$ en $\dfrac{3}{5} = \dfrac{x}{45}$:
    \[ 3\cdot 45 = 5x \;\Longrightarrow\; 135 = 5x
       \;\Longrightarrow\; x = \destacar{27} \]
    \azul{Verificación:} $\tfrac{27}{45} = \tfrac{3}{5}$ \checkmark
  \end{exampleblock}

  \vspace{0.06cm}
  {\tiny\color{uniGris} \destacar{Una razón no es una fracción del total.}
  ``3 a 2'' no significa que las mujeres sean $\tfrac{3}{2}$ del grupo, sino
  $\tfrac{3}{5}$ de él. Confundir las dos cosas es el error clásico del tema.}
\end{frame}

\begin{frame}{Regla de tres: directa e inversa}
  \scriptsize
  \begin{columns}[T,onlytextwidth]
    \begin{column}{0.48\textwidth}
      \begin{block}{Directa: más da más}
        \tiny
        $8$ kg de un insumo cuestan $\$36\,000$. ¿Cuánto cuestan $14$ kg?
        \[ \frac{8}{36\,000} = \frac{14}{x}
           \;\Longrightarrow\; x = \frac{36\,000\cdot 14}{8} \]
        \[ x = \destacar{\$63\,000} \]
      \end{block}
    \end{column}
    \begin{column}{0.48\textwidth}
      \begin{alertblock}{Inversa: más da menos}
        \tiny
        $6$ operarios tardan $15$ días. ¿Cuánto tardan $10$?
        \[ 6\cdot 15 = 10\cdot x \;\Longrightarrow\;
           x = \frac{90}{10} \]
        \[ x = \destacar{9\ \text{días}} \]
      \end{alertblock}
    \end{column}
  \end{columns}

  \vspace{0.12cm}
  \begin{block}{Cómo decidir cuál es}
    \tiny
    Pregúntese: \destacar{si una magnitud aumenta, ¿la otra aumenta o
    disminuye?} Aumenta $\Rightarrow$ \azul{directa} (se multiplica en cruz).
    Disminuye $\Rightarrow$ \azul{inversa} (el producto de cada pareja es
    constante).
  \end{block}

  \vspace{0.06cm}
  {\tiny\color{uniGris} \destacar{Control de sentido común:} si eran $6$
  operarios en $15$ días y ahora son más operarios, la respuesta \azul{tiene}
  que ser menos de $15$ días. Si le da más, aplicó la regla equivocada.}
\end{frame}

\begin{frame}{Porcentajes: los cuatro tipos de pregunta}
  \scriptsize
  \begin{center}
  \begin{tabular}{@{}p{3.6cm}p{4.0cm}p{4.2cm}@{}}
    \toprule
    \textbf{Pregunta} & \textbf{Se resuelve} & \textbf{Ejemplo} \\
    \midrule
    ¿Cuánto es el $p\%$ de $N$? &
      $N\cdot\dfrac{p}{100}$ &
      $18\%$ de $2\,450 = \azul{441}$ \\
    \addlinespace
    ¿Qué porcentaje es $a$ de $N$? &
      $\dfrac{a}{N}\cdot 100$ &
      $63$ de $180 = \azul{35\%}$ \\
    \addlinespace
    Si $a$ es el $p\%$, ¿cuánto es el total? &
      $\dfrac{a\cdot 100}{p}$ &
      $125$ es el $25\% \Rightarrow \azul{500}$ \\
    \addlinespace
    ¿En qué porcentaje varió? &
      $\dfrac{\text{final}-\text{inicial}}{\text{inicial}}\cdot 100$ &
      De $45$ a $63$: $\azul{+40\%}$ \\
    \bottomrule
  \end{tabular}
  \end{center}

  \vspace{0.1cm}
  \begin{alertblock}{El denominador de una variación es siempre el valor inicial}
    \tiny
    De $80$ a $68$: la variación es $\tfrac{-12}{80}\cdot 100
    = \destacar{-15\%}$, no $\tfrac{-12}{68}\cdot 100 = -17{,}6\%$.
    \azul{Se compara contra el punto de partida}, no contra el de llegada.
  \end{alertblock}
\end{frame}

\begin{frame}{Los porcentajes no se suman}
  \scriptsize
  \begin{columns}[T,onlytextwidth]
    \begin{column}{0.48\textwidth}
      \begin{alertblock}{Dos descuentos sucesivos}
        \tiny
        Un $20\%$ y luego un $10\%$ sobre un precio de $100$:
        \[ 100\cdot 0{,}80\cdot 0{,}90 = 72 \]
        El descuento total es del \destacar{28\%}, no del $30\%$.
        El segundo se aplica sobre lo ya rebajado.
      \end{alertblock}
    \end{column}
    \begin{column}{0.48\textwidth}
      \begin{alertblock}{Subir y bajar lo mismo}
        \tiny
        Un precio sube $20\%$ y luego baja $20\%$:
        \[ 100\cdot 1{,}20\cdot 0{,}80 = 96 \]
        \destacar{No se vuelve al precio original.} Se pierde un $4\%$,
        porque la bajada se calcula sobre una base mayor.
      \end{alertblock}
    \end{column}
  \end{columns}

  \vspace{0.12cm}
  \begin{block}{El truco del factor}
    \tiny
    Aumentar un $p\%$ es multiplicar por $\left(1+\tfrac{p}{100}\right)$;
    disminuir, por $\left(1-\tfrac{p}{100}\right)$. Encadenar variaciones es
    \destacar{multiplicar factores}, nunca sumar porcentajes.
    \[ \text{Tres aumentos del } 10\%:\quad
       100\cdot 1{,}1^{3} = \destacar{133{,}10} \quad
       \text{(no } 130\text{)} \]
  \end{block}

  \vspace{0.06cm}
  {\tiny\color{uniGris} \azul{Para quitar el IVA:} si un precio con IVA del
  $19\%$ es $\$84\,000$, la base es $84\,000\div 1{,}19 \approx
  \$70\,588$. \destacar{No} se le resta el $19\%$ a $84\,000$.}
\end{frame}

\begin{frame}{Repartos proporcionales}
  \scriptsize
  \begin{exampleblock}{Repartir $\$4\,200$ en la razón $2:3:5$}
    \tiny
    \textbf{1.} Sume las partes: $2+3+5 = 10$. \quad
    \textbf{2.} Halle el valor de una parte: $4\,200\div 10 = 420$. \quad
    \textbf{3.} Multiplique:
    \[ 2\cdot 420 = \destacar{840} \qquad
       3\cdot 420 = \destacar{1\,260} \qquad
       5\cdot 420 = \destacar{2\,100} \]
    \azul{Verificación:} $840+1\,260+2\,100 = 4\,200$ \checkmark
  \end{exampleblock}

  \vspace{0.1cm}
  \begin{block}{Otro caso}
    \tiny
    Repartir $91$ en la razón $3:4$: las partes suman $7$, cada parte vale
    $13$, y el reparto es $\destacar{39}$ y $\destacar{52}$.
  \end{block}

  \vspace{0.08cm}
  {\tiny\color{uniGris} \destacar{La verificación es obligatoria y gratis:}
  las partes repartidas siempre tienen que sumar el total. Si no suman, hubo
  un error de división. Es el único ejercicio del examen que puede
  comprobarse por completo en diez segundos.}
\end{frame}

\begin{frame}{Ejercicio 1 — proporciones y porcentajes}
  \scriptsize
  \begin{alertblock}{En parejas, 12 minutos}
    \tiny
    \textbf{1.} Halle $x$: \ $\tfrac{7}{12}=\tfrac{21}{x}$ \qquad
    \textbf{2.} $35\%$ de $640$ \qquad
    \textbf{3.} ¿Qué porcentaje es $63$ de $180$? \qquad
    \textbf{4.} Reparta $60$ en la razón $1:2:3$ \qquad
    \textbf{5.} Un artículo pasa de $\$45\,000$ a $\$63\,000$: ¿cuánto subió
    en porcentaje?
  \end{alertblock}

  \vspace{0.05cm}
  \begin{block}{Solución}
    \tiny
    \textbf{1.} $7x = 12\cdot 21 = 252 \Rightarrow x=\destacar{36}$ \qquad
    \textbf{2.} $640\cdot 0{,}35 = \destacar{224}$ \qquad
    \textbf{3.} $\tfrac{63}{180}\cdot 100 = \destacar{35\%}$

    \vspace{0.08cm}
    \textbf{4.} Partes: $1+2+3=6$; cada parte $=10$:
    $\destacar{10,\ 20,\ 30}$ \qquad
    \textbf{5.} $\tfrac{63\,000-45\,000}{45\,000}\cdot 100
    = \tfrac{18}{45}\cdot 100 = \destacar{+40\%}$
  \end{block}
\end{frame}

% ===========================================================================
\seccion{Repaso integrador de la Unidad 1}
% ===========================================================================

\begin{frame}{Mapa de la Unidad 1}
  \scriptsize
  \begin{center}
  \begin{tabular}{@{}clp{7.6cm}@{}}
    \toprule
    \textbf{Tema} & \textbf{Sesión} & \textbf{Lo que se evalúa} \\
    \midrule
    1.1 & 1 &
      Clasificar en $\mathbb{N},\mathbb{Z},\mathbb{Q},\mathbb{I}$; decimal a
      fracción; propiedades; orden; intervalos; valor absoluto \\
    \addlinespace
    1.2 & 2 &
      Ley de los signos; potencias de enteros; jerarquía de operaciones;
      factorización, MCD y mcm \\
    \addlinespace
    1.3 & 3 &
      Simplificar; sumar, restar, multiplicar y dividir fracciones;
      fracciones complejas; jerarquía \\
    \addlinespace
    1.4 & 4 &
      Extensión y comprensión; $\in$ y $\subseteq$; unión, intersección,
      diferencia, complemento; Venn y conteo \\
    \addlinespace
    1.5 & 5 &
      Razones, proporciones, regla de tres, porcentajes, repartos \\
    \bottomrule
  \end{tabular}
  \end{center}

  \vspace{0.08cm}
  {\tiny\color{uniGris} El Examen 1 evalúa \destacar{los cinco temas}, no una
  selección. Si algo quedó flojo, esta semana es el momento: la Unidad 2
  arranca de cero en cuanto a temas, pero da por sabido todo lo de arriba.}
\end{frame}

\begin{frame}{Simulacro del Examen 1 — parte A}
  \begin{alertblock}{Individual, 20 minutos, sin apuntes}
    \scriptsize
    \begin{enumerate}\setlength{\itemsep}{0.0cm}
      \item Clasifique en $\mathbb{N},\mathbb{Z},\mathbb{Q},\mathbb{I}$: \
            $\sqrt{121}$, \ $\sqrt{28}$, \ $-\tfrac{18}{6}$, \
            $0{,}\overline{36}$
      \item Convierta a fracción: \ $0{,}\overline{36}$ \ y \
            $1{,}4\overline{3}$
      \item Calcule: \ $-4^{2}+(-3)^{3}-20\div(-4)\cdot 3$
      \item Calcule: \ $\dfrac{5}{8}+\dfrac{2}{3}-\dfrac{1}{4}$
      \item Calcule: \ $\dfrac{2}{5}\div\dfrac{3}{10}$
      \item Halle $\mathrm{MCD}(126,210)$ y $\mathrm{mcm}(126,210)$
    \end{enumerate}
  \end{alertblock}
\end{frame}

\begin{frame}{Simulacro — solución de la parte A}
  \scriptsize
  \begin{columns}[T,onlytextwidth]
    \begin{column}{0.49\textwidth}
      \begin{block}{Puntos 1 a 3}
        \tiny
        \textbf{1.} $\sqrt{121}=11$: $\mathbb{N},\mathbb{Z},\mathbb{Q}$. \
        $\sqrt{28}=2\sqrt{7}$: $\mathbb{I}$. \
        $-\tfrac{18}{6}=-3$: $\mathbb{Z},\mathbb{Q}$. \
        $0{,}\overline{36}$: $\mathbb{Q}$.

        \vspace{0.08cm}
        \textbf{2.} $99x=36 \Rightarrow x=\destacar{\tfrac{4}{11}}$; \quad
        $100x-10x = 143{,}\overline{3}-14{,}\overline{3} = 129
        \Rightarrow x=\tfrac{129}{90}=\destacar{\tfrac{43}{30}}$

        \vspace{0.08cm}
        \textbf{3.} $-16 + (-27) - (-5)(3) = -16-27+15
        = \destacar{-28}$
      \end{block}
    \end{column}
    \begin{column}{0.49\textwidth}
      \begin{block}{Puntos 4 a 6}
        \tiny
        \textbf{4.} m.c.d. $=24$:
        $\tfrac{15+16-6}{24} = \destacar{\tfrac{25}{24}}$

        \vspace{0.1cm}
        \textbf{5.} $\tfrac{2}{5}\cdot\tfrac{10}{3}
        = \tfrac{20}{15} = \destacar{\tfrac{4}{3}}$

        \vspace{0.1cm}
        \textbf{6.} $126=2\cdot 3^{2}\cdot 7$, \
        $210 = 2\cdot 3\cdot 5\cdot 7$
        \[ \mathrm{MCD}=\destacar{42} \qquad
           \mathrm{mcm}=\destacar{630} \]
      \end{block}
    \end{column}
  \end{columns}

  \vspace{0.06cm}
  {\tiny\color{uniGris} En el \destacar{3}, $20\div(-4)\cdot 3 = -5\cdot 3
  = -15$, y como va restando queda $+15$. Verificación del \destacar{6}:
  $42\cdot 630 = 26\,460 = 126\cdot 210$ \checkmark}
\end{frame}

\begin{frame}{Simulacro del Examen 1 — parte B}
  \begin{alertblock}{Individual, 20 minutos}
    \scriptsize
    \begin{enumerate}\setlength{\itemsep}{0.0cm}
      \setcounter{enumi}{6}
      \item Calcule: \ $|6-14| - |-5|$
      \item Escriba en notación de intervalo: ``desde $-3$ inclusive hasta
            $8$ sin incluir''.
      \item Con $U=\{1,\dots,10\}$, \ $A$ los pares y $B$ los múltiplos de
            $3$: halle $A\cup B$, $A\cap B$ y $\#(A\cup B)$.
      \item De $70$ usuarios, $40$ usan el canal virtual, $35$ el
            presencial y $18$ los dos. ¿Cuántos no usan ninguno?
      \item Halle $x$: \ $\dfrac{x}{9}=\dfrac{12}{27}$
      \item Un precio de $\$120\,000$ recibe un descuento del $15\%$ y luego
            otro del $10\%$. ¿Cuál es el precio final y cuál el descuento
            total en porcentaje?
    \end{enumerate}
  \end{alertblock}
\end{frame}

\begin{frame}{Simulacro — solución de la parte B}
  \scriptsize
  \begin{columns}[T,onlytextwidth]
    \begin{column}{0.49\textwidth}
      \begin{block}{Puntos 7 a 9}
        \tiny
        \textbf{7.} $8-5 = \destacar{3}$

        \vspace{0.08cm}
        \textbf{8.} $\destacar{[-3,\ 8)}$

        \vspace{0.08cm}
        \textbf{9.} $A=\{2,4,6,8,10\}$, $B=\{3,6,9\}$
        \[ A\cup B = \{2,3,4,6,8,9,10\},\quad A\cap B=\{6\} \]
        $\#(A\cup B) = 5+3-1 = \destacar{7}$
      \end{block}
    \end{column}
    \begin{column}{0.49\textwidth}
      \begin{block}{Puntos 10 a 12}
        \tiny
        \textbf{10.} $40+35-18 = 57$; \
        $70-57 = \destacar{13}$

        \vspace{0.08cm}
        \textbf{11.} $27x = 108 \Rightarrow x=\destacar{4}$

        \vspace{0.08cm}
        \textbf{12.} $120\,000\cdot 0{,}85\cdot 0{,}90
        = \destacar{\$91\,800}$; \
        descuento total $= 23{,}5\%$ \azul{(no 25\%)}
      \end{block}
    \end{column}
  \end{columns}

  \vspace{0.06cm}
  {\tiny\color{uniGris} El punto \destacar{12} es el que separa: $0{,}85\cdot
  0{,}90 = 0{,}765$, es decir se paga el $76{,}5\%$ y se descuenta el
  $23{,}5\%$. \azul{Los porcentajes encadenados se multiplican, no se
  suman.}}
\end{frame}

\begin{frame}{Errores frecuentes de toda la Unidad 1}
  \scriptsize
  \begin{center}
  \begin{tabular}{@{}p{5.4cm}p{6.6cm}@{}}
    \toprule
    \textbf{Error} & \textbf{Cómo evitarlo} \\
    \midrule
    $-5 > -2$ &
      Con negativos el orden se invierte \\
    \addlinespace
    $-2^{4}=16$ &
      Sin paréntesis el menos no se eleva: $-16$ \\
    \addlinespace
    $\frac{3}{4}+\frac{5}{6}=\frac{8}{10}$ &
      Iguale denominadores; nunca los sume \\
    \addlinespace
    Invertir la primera fracción al dividir &
      Se invierte la segunda \\
    \addlinespace
    $\#(A\cup B)=\#A+\#B$ &
      Reste $\#(A\cap B)$ salvo que sean disjuntos \\
    \addlinespace
    Sumar porcentajes sucesivos &
      Multiplique factores: $0{,}85\cdot 0{,}90$ \\
    \addlinespace
    Usar el valor final como denominador &
      La variación se calcula sobre el valor \azul{inicial} \\
    \bottomrule
  \end{tabular}
  \end{center}
\end{frame}

\begin{frame}{Cierre y preparación del Examen 1}
  \scriptsize
  \begin{alertblock}{Examen 1 (35\%) — semana 6, 14 al 19 de septiembre}
    \tiny
    Individual, en \azul{Aulas Virtuales (Moodle)}, presencial en el aula, con
    tiempo límite. Evalúa los cinco temas de la Unidad 1.
  \end{alertblock}

  \vspace{0.1cm}
  \begin{block}{Cómo prepararlo esta semana}
    \tiny
    \textbf{1.} Rehaga el simulacro de hoy \azul{cronometrado} y sin mirar las
    soluciones. \quad
    \textbf{2.} Repase los ejercicios propuestos de las sesiones 1 a 5; las
    respuestas están en las diapositivas. \quad
    \textbf{3.} Resuelva la guía de la unidad en Moodle. \quad
    \textbf{4.} Lleve calculadora, aunque casi todo se puede hacer a mano.
  \end{block}

  \vspace{0.1cm}
  \begin{exampleblock}{Lo que viene después}
    \tiny
    En la sesión 7 empieza la \azul{Unidad 2}: potenciación, radicación,
    racionalización y logaritmos. Es la unidad más ``mecánica'' del curso —hay
    reglas claras que se aplican— y por eso suele subir el promedio de quien
    practica.
  \end{exampleblock}

  \vspace{0.06cm}
  {\tiny\color{uniGris} Consultas y asesorías: hasta el día anterior al
  examen, por correo o al final de la clase.}
\end{frame}

\begin{frame}{Ejercicios propuestos}
  \vspace{-0.1cm}
  \begin{beamercolorbox}[wd=\textwidth,sep=0.10cm]{block body}
    \centering\scriptsize
    \textbf{1--6:} proporciones y regla de tres. \quad
    \textbf{7--14:} porcentajes. \quad
    \textbf{15--20:} repaso de toda la unidad.
  \end{beamercolorbox}
  \listatiny
  \vspace{0.15cm}
  \begin{columns}[T,onlytextwidth]
    \begin{column}{0.47\textwidth}
      \begin{enumerate}\setlength{\itemsep}{0.16cm}
        \item $\tfrac{x}{8}=\tfrac{15}{24}$
        \item $\tfrac{5}{x}=\tfrac{20}{36}$
        \item $12$ kg cuestan $\$54\,000$: ¿$20$ kg?
        \item $4$ máquinas tardan $18$ h: ¿$6$ máquinas?
        \item Reparta \$3\,600 en razón $1:3:5$
        \item Reparta $84$ en razón $2:5$
        \item $24\%$ de $1\,750$
        \item ¿Qué \% es $56$ de $224$?
        \item $90$ es el $15\%$ de\dots
        \item De $250$ a $310$: variación \%
      \end{enumerate}
    \end{column}
    \begin{column}{0.50\textwidth}
      \begin{enumerate}\setlength{\itemsep}{0.16cm}
        \setcounter{enumi}{10}
        \item De $60$ a $48$: variación \%
        \item Precio $\$95\,200$ con IVA $19\%$: base
        \item Descuentos del $30\%$ y luego $20\%$: \% total
        \item Sube $25\%$ y luego baja $25\%$: \% neto
        \item Convierta $0{,}\overline{54}$ a fracción
        \item Calcule $-3^{2}+(-4)^{2}-18\div(-3)$
        \item Calcule $\tfrac{3}{4}-\tfrac{5}{6}\cdot\tfrac{3}{10}$
        \item $\mathrm{mcm}(15,20,24)$
        \item Con $U=\{1,\dots,12\}$, $A$ múltiplos de $4$, $B$ múltiplos
              de $6$: halle $A\cap B$
        \item Calcule $|3-11|+|-4|-|2-2|$
      \end{enumerate}
    \end{column}
  \end{columns}
\end{frame}

\begin{frame}{Respuestas}{Autoevaluación de los ejercicios propuestos}
  \listatiny
  \vspace{0.15cm}
  \begin{columns}[T,onlytextwidth]
    \begin{column}{0.47\textwidth}
      \begin{enumerate}\setlength{\itemsep}{0.16cm}
        \item $x=5$
        \item $x=9$
        \item $\$90\,000$
        \item $12$ h
        \item $400,\ 1\,200,\ 2\,000$
        \item $24$ y $60$
        \item $420$
        \item $25\%$
        \item $600$
        \item $+24\%$
      \end{enumerate}
    \end{column}
    \begin{column}{0.50\textwidth}
      \begin{enumerate}\setlength{\itemsep}{0.16cm}
        \setcounter{enumi}{10}
        \item $-20\%$
        \item $\$80\,000$
        \item $44\%$ \ {\color{uniGris}($0{,}7\cdot 0{,}8=0{,}56$)}
        \item $-6{,}25\%$
        \item $\tfrac{6}{11}$
        \item $13$
        \item $\tfrac{1}{2}$
        \item $120$
        \item $\{12\}$
        \item $12$
      \end{enumerate}
    \end{column}
  \end{columns}
\end{frame}

\begin{frame}{Referencias}
  \footnotesize
  \begin{itemize}
    \item[{[1]}] Ramos del Olmo, Á. M. (2016). \textit{Matemáticas básicas
      para el acceso a la universidad}. Ediciones Pirámide.
      \texttt{https://elibro.net/es/ereader/uniminuto/49134}
    \item[{[2]}] González Pulido, J. W., Chacón Benavides, J. A. \&
      Fonseca Correa, L. Á. (2023). \textit{Matemática básica}
      (1.\textsuperscript{a} ed.). Universidad Mariana.
      \texttt{https://elibro.net/es/ereader/uniminuto/280647}
    \item[{[3]}] Torrijos, M. (2018). \textit{Apuntes de matemática básica}.
      Universidad Católica de Colombia.
      \texttt{https://elibro.net/es/ereader/uniminuto/197048}
    \item[{[4]}] Ramírez, A. \& Rojas, L. (2016). \textit{Matemáticas básicas
      con aplicaciones a la ingeniería}. Ecoe Ediciones.
      \texttt{https://elibro.net/es/ereader/uniminuto/114355}
  \end{itemize}

  \vspace{0.3cm}
  \begin{center}
    {\Large\color{uniAzul}\textbf{¿Preguntas?}}
  \end{center}
\end{frame}

\end{document}
